\section{The terminal Wilson family}
\label{sec:wilson-startup-construction}
The shooting argument needs a continuous initial family on the entire
terminal coupling band, with one fixed local dictionary and common
summable bounds. We construct it directly from the Wilson density.
This also explains why its terminal path can be chosen before any
trajectory is selected.

Fix a compact gauge group as in the construction, a clipping profile
with direct coefficient $0<c_0\le1$, linear coefficient
$0<c_1<c_0/16$, so that in particular $c_1\le1/4$, and relaxation
length $\lambda>0$. Write $q=1-c_0$.
For a plaquette holonomy $u$, put
\[
 r(u)=\operatorname{chord}(u),\qquad
 k(u)=\operatorname{curv}(u),\qquad
 W(u)=\tfrac12r(u)^2.
\]
Thus $0\le W\le2$ and $\|k(u)\|\le r(u)$. The gauge normalization
gives $k(\exp X)=X+O(\|X\|^2)$ near zero.

\subsection{A curvature-coordinate remainder}
The analytic inverse function theorem applied to $X\mapsto k(\exp X)$
gives a local inverse $v$ near zero. Set $u(x)=\exp v(x)$ and
$w(x)=W(u(x))$. Conjugation covariance of curvature and uniqueness
of the local inverse show
$w(\operatorname{Ad}_g x)=w(x)$ as germs at zero. The identities
$k(u^{-1})=-k(u)$ and $W(u^{-1})=W(u)$ give $w(-x)=w(x)$.
The first and second derivatives at zero are respectively zero and
the Euclidean inner product. Thus the fourth Taylor polynomial is
\[
 T_4w(x)=\tfrac12\|x\|^2+Q_4(x),\qquad
 Q_4(x)=\frac1{4!}D^4w(0)[x,x,x,x].
\]
The odd terms vanish by evenness. The quartic form $Q_4$ is
adjoint-invariant, by differentiation of the conjugation identity.
Define the global smooth residual
\[
 h(u)=W(u)-\tfrac12\|k(u)\|^2-Q_4(k(u)).
\]
For every smooth family through the identity its composition with $h$
has vanishing derivatives through order four. Indeed near the
identity $u=u(k(u))$, so this is the Taylor remainder
$(w-T_4w)(k(u))$; the chain rule preserves the stated vanishing.

We need more than a Taylor estimate at the all-identity link field.
For every derivative order $K$ and spatial power $P$ there are
$C>0$ and $t_0\in(0,1]$, independent of the torus, such that
\begin{equation}\label{eq:startup-residual-price}
 \bigl|\partial_{a_1}\cdots\partial_{a_l}
              h(U_0U_1U_2^{-1}U_3^{-1})\bigr|
       \le Ct^5,\qquad 0\le l\le K,
\end{equation}
when the plaquette chord is at most $t\le t_0$ and the four-link
directions have norm at most $t5^P$. Derivatives here are successive
left increments on the links.

Here is a uniform proof. For fixed links $U_1,U_2,U_3$ and a small
face logarithm $x$, reconstruct the first link as
\[
             U_0(x)=\exp(x)U_3U_2U_1^{-1}.
\]
The face holonomy is then $\exp x$. After $l$ left differentiations
the resulting smooth function of $x$ vanishes to order
$(5-l)_+=\max(5-l,0)$ at zero. Each differentiation can lower the
vanishing order by at most one: introduce its increment parameter
as one further variable, differentiate the vanishing Taylor
polynomial, and restrict that parameter to zero. All remaining
links range over a compact group and the directions, after division
by $t$, range over a fixed compact ball. Taylor's integral remainder
therefore gives a common bound $C_l\|x\|^{(5-l)_+}$ for the
unscaled derivatives. A fixed local logarithm satisfies
$\|x\|\le C_{\log}r(u)$ when $r(u)$ is sufficiently small.
Restoring the $l$ factors of $t$ gives
$C_l'C_{\log}^{(5-l)_+}t^{l+(5-l)_+}\le C_l''t^5$,
since $t\le1$. Take the largest constant and smallest threshold
over $0\le l\le K$. For a real lattice plaquette, the four links
may coincide because of wrapping; substitution into the four-link
identity remains valid. Each lies within distance four of any
chosen plaquette corner, so the normalized spatially weighted
directions obey the required bound $t5^P$. This proves
\eqref{eq:startup-residual-price} on every healthy plaquette,
including equality in the health threshold.

\subsection{The principal path and the nonnegative debt}
Choose the terminal path to have one row
$\sqrt q\,k(u)$ at each elementary plaquette, quartic term $qQ_4$,
and zero cubic and loop terms. Its zero-momentum metric is
$c_0I+qI=I$. Each row slot is an elementary plaquette slot with empty
walk, hence of length four in the path-length convention of
Section~\ref{subsec:path-stationary-chart}. Each quartic atom has four
such legs, all at the same plaquette, and therefore has length
sixteen although its geometric support is a single plaquette. Its only
nonzero quartic coefficients form a finite list. Hence at any fixed
positive strip width $\sigma$ the weighted path sums are finite: the
row sum is $\sqrt q\,e^{4\sigma}$ and the quartic sum is $e^{16\sigma}$
times the finite sum of the absolute quartic coefficients. Choosing
the stage path allowance above their maximum, and the column
allowance above the prescribed factorial, supplies the startup
profile before choosing block lengths. Its uncut principal action
differs from the Wilson action by $\sum qh(u)$, so the two actions
have identical jets through order four at the identity.

Fix an integer support scale $B\ge2$ and threshold $t>0$ satisfying
\begin{equation}\label{eq:startup-relaxation}
                  2\le t e^{(B-2)/\lambda}.
\end{equation}
At a plaquette with base site $x$, the row is kept when every
plaquette based at $y$ has chord at most $te^{d(x,y)/\lambda}$.
Its base action is
\[
 b(u)=
 \begin{cases}c_0r(u)^2/2,&r(u)\le t,\\c_1tr(u),&r(u)>t,\end{cases}
       +\mathbf1_{\rm kept}\,q\|k(u)\|^2/2.
\]
The quartic term is masked by health of the full four-corner
plaquette support. Define the debt at that plaquette by
\[
 d(U)=
 \begin{cases}
 W(u)-c_1tr(u),&r(u)>t,\\
 \mathbf1_{\rm healthy}\,q\|k(u)\|^2/2+
 \mathbf1_{\rm not\ healthy}\,qW(u)
       -\mathbf1_{\rm kept}\,q\|k(u)\|^2/2,&r(u)\le t .
 \end{cases}
\]
A direct split into the healthy, large-chord and remaining cases gives
\begin{equation}\label{eq:startup-face-split}
 W(u)=b(u)+q\mathbf1_{\rm healthy}Q_4(k(u))
             +q\mathbf1_{\rm healthy}h(u)+d(U).
\end{equation}
The debt satisfies $0\le d(U)\le2$, since
$q\|k(u)\|^2/2\le qW(u)\le2$. On a large-chord plaquette,
\[
 d(U)\ge r(u)^2/2-tr(u)/4\ge t^2/4.
\]
If the base site is farther than $B$ from every flag, the plaquette
is healthy and the row is kept. For nearby plaquettes this follows
from absence of flags; for the remaining plaquettes
\eqref{eq:startup-relaxation} and the global chord bound $r\le2$
make the kept test automatic. Thus its debt vanishes there.
The same argument shows that debt at a base site can be determined
from the field in its radius-$B$ neighborhood, with the prescribed
boundary flag information.

Average \eqref{eq:startup-face-split} over signed coordinate changes.
Let $s_*=1/|\mathcal S|$, where $\mathcal S$ is that finite symmetry
group, and index the oriented plaquettes by $i=(s,x,b)$.
Their supports are the actual transformed four-corner supports
$Y_i$, with native site $x$, residual $h_i$ and debt $d_i$.
The signed average is exact, because the Wilson energy is invariant:
\begin{equation}\label{eq:startup-signed-split}
 \beta E_W(U)=A^{\rm top}_\beta(U)
       +\beta s_*q\sum_i\mathbf1_{\rm healthy(Y_i)}h_i(U)
       +\beta s_*\sum_i d_i(U).
\end{equation}
This uses the transformed supports and native sites in every channel.

\subsection{Forced packets and the exact Wilson density}
Let $\mathcal D(U)$ be the flag set. Join two flags when their closed
radius-$2B$ balls overlap, and let $D$ run over the connected
components for this relation. Put
\[
 P(D)=\bigcup_{x\in D}B_{2B}(x),\qquad
 C_D(U)=s_*\sum_{i:\,x_i\in B_B(D)}d_i(U).
\]
The regions $P(D)$ of distinct components are disjoint and
$P(D)\cap\mathcal D(U)=D$. Every nonzero debt lies in exactly one
of these sums: it lies within $B$ of a flag, while membership in
two such neighborhoods would join the two flag components.
Consequently
\[
                 \sum_D C_D(U)=s_*\sum_i d_i(U).
\]

The energy price is
\begin{equation}\label{eq:startup-packet-energy}
                    C_D(U)\ge |D|s_*t^2/16.
\end{equation}
To prove it, every flag is a corner of a large-chord plaquette.
Such a plaquette has at most four corners and its base is within
distance two of the flag, hence in $B_B(D)$. If $\mathcal A$ is
the set of these plaquettes in the identity channel, then
$|D|\le4|\mathcal A|$. Each contributes at least $t^2/4$;
all other debts are nonnegative. Multiplication by $s_*$ proves
\eqref{eq:startup-packet-energy}.

In a fixed dictionary include every nonempty connected flag set $D$
with address $(P(D),D)$, and put
\[
 K_D(\beta,U)=
 \mathbf1_{\{P(D)\cap\mathcal D(U)=D\}}e^{-\beta C_D(U)}.
\]
Other addresses have zero coefficient. The forced gas sums products
over disjoint supports with total flag inventory $\mathcal D(U)$.
Exactly one family can have nonzero product: the components of
$\mathcal D(U)$ just constructed. Indeed disjoint connected regions
covering that inventory give its component decomposition. Thus the
gas equals $e^{-\beta s_*\sum_i d_i(U)}$, including the empty
product when there are no flags. Put the regular residuals
$\beta s_*q h_i$ in their healthy supports and take every extra
logarithmic attachment to be zero. Equation
\eqref{eq:startup-signed-split} now proves that the full density is
exactly $e^{-\beta E_W(U)}$ at every field.

All coefficients are gauge-invariant and measurable; the local
residuals are smooth in left increments. Their value and Hessian
at the identity vanish. Permuting the signed channels and lattice
sites gives the required translation and signed-coordinate
covariance. The packet charge reads only its support with the given
flags, since the radius-$B$ neighborhood of a site in $B_B(D)$
lies in $P(D)$. These are the locality conditions of the startup
family.

\subsection{Uniform prices on the whole coupling band}
Let $a>0$ be the support exponent. A plaquette support has cover
number at most $5^4$. At a fixed native site there are
$|\mathcal S|$ times the number $n_\square$ of plaquette directions
regular entries. By \eqref{eq:startup-residual-price} their regular
budget for $0\le\beta\le v$ is at most
\begin{equation}\label{eq:startup-regular-book}
             |\mathcal S|n_\square e^{a5^4}\,vCt^5.
\end{equation}
The constants depend on the stage, not on volume or cutoff.

Choose positive numbers $\sigma,\kappa,\gamma$ and a lower band
endpoint $l$ such that
\begin{align}
 \kappa+\sigma&\le l\,s_*t^2/16,\label{eq:startup-loan}\\
 (4B+1)^4e^{(a+2)5^4-\sigma}&\le\gamma,\qquad
 (2B+1)^4\gamma\le1,\label{eq:startup-single-hit}\\
 (2B+1)^4(e^\gamma-1)&\le1.\label{eq:startup-margin}
\end{align}
For $\beta\ge l$, \eqref{eq:startup-packet-energy} implies
\[
               |K_D(\beta,U)|\le e^{-\kappa}e^{-\sigma|D|}.
\]
The elementary ball supports $B_{2B}(x)$ have cover number at most
$5^4$, and at most $(4B+1)^4$ of them contain a specified site.
Thus \eqref{eq:startup-single-hit} bounds their rooted support sum
with exponent $a+2$ by $\gamma$. Apply the connected-support
estimate \eqref{eq:ope-forest-price}, with spare exponent one, to
these ball supports. It gives the packet majorant
\[
 \sup_x\sum_{D:\,x\in P(D)}
          e^{-\sigma|D|}e^{(a+1)\operatorname{cov}_B(P(D))}
                        \le e^\gamma-1.
\]
This is one common summable majorant for the entire band
$\beta\ge l$. The regular and forced budgets, and the extra
one-unit support reserve, have now been proved. After the
finite gas grouping the forced cost is at most
$e^{-\kappa}(e^{e^\gamma-1}-1)$, by the same connected-support
estimate with \eqref{eq:startup-margin}.

We verify that the actual schedule pays these inequalities. Write
$s=\log\widehat\beta$ and set
\[
 M=e^{(\log s)^2},\quad t=e^{-s/2}M,\quad B=\lceil s^2\rceil,
 \quad \gamma=e^{-M/2},\quad
 \delta_E=e^{-11s/10},\quad \delta_K=e^{-M^{1/5}}.
\]
At a fixed stage take $\sigma=M$, $l=e^s/2$ and $\kappa=M$.
Then $lt^2s_*/16=s_*M^2/32\ge2M$ eventually, giving
\eqref{eq:startup-loan}. Since $B$ grows polynomially in $s$
while $M$ grows faster than every such polynomial,
\eqref{eq:startup-single-hit} and \eqref{eq:startup-margin}
hold eventually. More explicitly $(4B+1)^4\le M/4$ and
$(a+2)5^4\le M/4$ eventually, so their product with
$e^{-M}$ is at most $e^{-M/2}$; also
$e^\gamma-1\le2\gamma$ once $\gamma\le1$.
Since $\gamma\to0$, we may also require
$e^{e^\gamma-1}-1\le1$. The grouped forced cost is then at most
$e^{-M}$ and hence at most $\delta_K$.

The logarithm of \eqref{eq:startup-regular-book}, with
$v=2e^s$, equals
\[ C-\frac32s+5(\log s)^2, \]
where $C$ is a stage constant.
It is at most $-11s/10$ for large $s$, because
$(\log s)^2=o(s)$. Also $t\to0$, $B\to\infty$, and
$\log(t e^{(B-2)/\lambda})=-s/2+(\log s)^2+
(B-2)/\lambda\to\infty$. Thus the residual threshold and
\eqref{eq:startup-relaxation} are paid at the same entry delay.

The stronger startup reserve needed by the later common services
is obtained without changing this family. For any fixed additional
support exponent $\eta\ge0$, replace $a$ by $a+\eta$ in the
preceding prices and delay the stage further. Take instead
$l=e^s/4$ and $\kappa=s_*M^2/128$. Eventually $M\le\kappa$,
and
\[
       \kappa+M\le2\kappa=s_*M^2/64=lt^2s_*/16.
\]
The same packet bounds then give a forced allowance at most
$e^{-s_*M^2/128}$ before grouping, with a still smaller
grouped cost after the reserve is paid. These choices depend on
the stage and its fixed extra requests, and work for all periods
and cutoffs realizing that stage. Taking their entry thresholds
together with the other already chosen thresholds therefore
preserves the order of choices in the common construction.

\subsection{Continuity and the initial history}
The dictionary and the masks above do not depend on $\beta$.
Regular coefficients depend linearly on $\beta$, so their
normalized derivative differences have the bound
$C|\beta-\beta'|t^5$. For a packet, $0\le C_D(U)\le
2s_*|\mathcal S|\,|\Lambda|n_\square$, and, for
$\beta,\beta'\ge0$,
\[
 |e^{-\beta C_D(U)}-e^{-\beta'C_D(U)}|
                         \le C_D(U)|\beta-\beta'|.
\]
There are finitely many nonzero addresses on a fixed torus.
Their weighted difference sums consequently tend to zero with
$\beta'-\beta$, uniformly in $U$. Together with the band-wide
majorants above this proves the required curve property and
invariant-atom continuity. No continuity of a health indicator
as a function of $U$ is asserted.

Pack the dictionary into the fixed terminal dictionary by its
coupling-independent injection. Zero unused entries contribute
neither cost nor logarithm, and every nonzero coefficient, support,
inventory, and density is preserved. The finite gas grouping is the
absorption of Section~\ref{sec:soft-log-fusion} applied to the packed
packet activities with all logarithms equal to zero. It is applied to
this same dictionary and these same coefficients, and it leaves the
regular entries unchanged. Its continuity in $\beta$ is
Lemma~\ref{lem:reset-absorption-continuity}: the first hypothesis is
the fixed-address continuity just proved, and the second holds with
the band-wide majorants $e^{-\kappa}e^{-\sigma|D|}$, of which only
finitely many are nonzero on the fixed torus. Together with the
grouped forced cost computed above, this gives a continuous family
on $[\widehat\beta/2,2\widehat\beta]$ with exact Wilson density,
regular allowance $\delta_E$ and grouped forced allowance
$\delta_K$. Its scalar constant is zero and its coupling is
the input parameter $\beta$.

We collect the result, separating the numerical conditions from the
two identifications with the history.

\begin{proposition}[The terminal family]
\label{prop:wilson-startup-family}
Let $N$ be a level of the scale schedule and let its stage have clip
$(c_0,c_1,\lambda)$ with $c_0\le1$, derivative order $K$, spatial
power $P$ and support exponent $a$. Let $t$, $B$ and
$\widehat\beta$ be the threshold, radius and reference coupling of
level $N$, and $\delta_E,\delta_K$ its regular and forced allowances.
Assume:
\begin{enumerate}
\item $B\ge2$ and the relaxation condition
 \eqref{eq:startup-relaxation};
\item $t\le t_0$, where $t_0$ is the threshold of
 \eqref{eq:startup-residual-price} for $K,P$;
\item the regular book \eqref{eq:startup-regular-book} with
 $v=2\widehat\beta$ is at most $\delta_E$;
\item there are $\sigma,\kappa,\gamma$ with $\kappa\ge0$ satisfying
 \eqref{eq:startup-loan}--\eqref{eq:startup-margin} for
 $l=\widehat\beta/2$, and
 $e^{-\kappa}(e^{e^\gamma-1}-1)\le\delta_K$.
\end{enumerate}
Then for every period there is a family $F_\beta$,
$\beta\in[\widehat\beta/2,2\widehat\beta]$, of representations at
level $N$ in one fixed dictionary, such that: the path of $F_\beta$ is
the startup path of this section in the stage clip; its coupling is
$\beta$ and its vacuum constant is zero; its density is
$e^{-\beta E_W}$ at every field; its regular and forced arrays have
the locality, symmetry and jet properties listed above, with regular
cost at most $\delta_E$ and forced cost at most $\delta_K$; and
$\beta\mapsto F_\beta$ is continuous in the labelled regular and
forced norms. Moreover the reserved labels are empty. The fixed
terminal dictionary contains, for every regular address and every
natural number, a reserved regular label with that address, and for
every forced support with nonempty inventory and every natural
number, a reserved forced label with that support and inventory;
these lists do not depend on $\beta$. The regular and forced address
lists of $F_\beta$ are those of the fixed dictionary, and the
regular, respectively forced, coefficient of $F_\beta$ on every
reserved label is zero.

Suppose in addition that the stationary plan of the history has the
startup path as its terminal path at $(N,N)$, as in
Proposition~\ref{prop:terminal-replacement}, and that the recipe of
the history is the default recipe built from that plan and schedule.
Then every $F_\beta$ satisfies the retained-history conditions at
level $N$ and the starter relation of that recipe.
\end{proposition}
\begin{proof}
The first part was proved in the preceding subsections: condition (1)
is used for the debt and its locality, (2) for the residual price,
(3) for the regular budget on $0\le\beta\le2\widehat\beta$, and (4)
for the packet majorant on $\beta\ge\widehat\beta/2$ and the grouped
forced cost. The statement on reserved labels is a property of the
fixed packing and grouping, valid for every dictionary and
coefficient family to which they are applied. The packing injection
takes its values among the unreserved labels, and every entry outside
its image is zero. A reserved forced label is the group consisting of
one reserved hard label alone. Every grouped label of the packed
family consists of packed labels, so a reserved forced label is not
among them, and every grouped entry outside the image of the packed
family is zero. For the second part, the retained-history conditions at
level $N$ are: the path is the terminal path of the plan, the clip
and budget are those of level $N$, the arrays satisfy the stage
bounds at the allowances $\delta_E,\delta_K$, the path satisfies its
stage bounds, and the representation is a start of the recipe. The
first is the assumed identity of paths; it is an equality of path
data, which agreement of jets through order four would not give. The
bounds are the first part and the path bounds of the plan. A start of
the default recipe is a representation with coupling in the terminal
band which is obtained by the fixed packing and grouping from some
dictionary and coefficient family carrying a pre-absorption
certificate at the stage allowances, whose raw density on the
terminal path of the plan, multiplied by the exponential of its
logarithmic corrections, is the Wilson density. Our pre-grouping
dictionary and coefficients are such a witness: the regular book and
the packet loan give the certificate, the logarithms are zero, and
the raw density identity is the one proved above, for the terminal
path by the assumed identity of paths. The default recipe imposes no
further permission condition on a start. The starter relation is the
same statement with the band condition recorded.
\end{proof}

The four numerical conditions hold on the actual schedule, at every
level $N$ including $N=0$, for the following reason. The preceding
subsection showed that for each fixed stage they hold for all
sufficiently large $s=\log\widehat\beta$, with $t,B,\delta_E,\delta_K$
the functions of $s$ displayed there. The resulting threshold depends
only on the stage, so it can be imposed before the schedule is
realized. The common selection must include it among the entry
requirements of that stage; this inclusion is made in
Section~\ref{sec:common-funded-history} and is an input here. Granting
it, a paid schedule enters each stage beyond its entry requirement
and has at every level the threshold, radius and allowances given by
those functions of its reference coupling, so the four conditions
hold at every level. No blocking step is involved in these
conditions, so the case $N=0$ needs no separate argument. The identification of
the terminal path is Proposition~\ref{prop:terminal-replacement}, and
the recipe is fixed in Section~\ref{sec:common-funded-history}. With
them the proposition gives the initial-family assertion used in
Section~\ref{sec:common-funded-history} for every $N$ and every
period.
